\documentclass[11pt,a4paper]{article}
\usepackage[utf8]{inputenc}
\usepackage[T1]{fontenc}
\usepackage[brazil]{babel}
\usepackage{geometry}
\usepackage{amsmath}
\geometry{margin=2.2cm}

\title{Braço robótico 2D (Pygame)}
\author{Computação Gráfica --- Aula 05}
\date{}

\begin{document}
\maketitle
\vspace{-1.2em}

\section*{Objetivo}
Simular visualmente um manipulador planar com dois elos (\emph{links}) de comprimentos fixos $L_1$ e $L_2$, articulado na base e no ``cotovelo'', usando a biblioteca \texttt{pygame}.

\section*{Cinemática direta}
A base desliza ao longo do eixo horizontal $x$ (altura fixa $y = C_y$). O primeiro elo sai da base com ângulo $\theta_1$ (graus) em relação ao eixo $x$. O segundo elo soma o ângulo relativo $\theta_2$ ao ângulo absoluto do primeiro segmento, de modo que a orientação do segundo elo é $\theta_1 + \theta_2$. A extremidade de um segmento de comprimento $L$ a partir de $(x_0,y_0)$ é
\begin{equation*}
  x = x_0 + L\cos(\theta),\qquad y = y_0 + L\sin(\theta),
\end{equation*}
com $\theta$ convertido de graus para radianos (\texttt{math.radians}). Assim obtém-se o ponto final do primeiro elo e, a partir dele, o da ponta do segundo.

\section*{Representação gráfica}
Desenham-se eixos cartesianos (referência), retângulo arredondado da base, segmentos espessos para os braços e círculos nas articulações e na ponta. O laço principal limpa a tela, redesenha e limita a taxa a 60 FPS.

\section*{Controles}
\texttt{A}/\texttt{D} ou setas esquerda/direita movem a base; \texttt{W}/\texttt{S} alteram $\theta_1$; setas cima/baixo alteram $\theta_2$. Incrementos de $4^\circ$ e repetição de tecla facilitam o ajuste contínuo. \texttt{Esc} encerra o programa.

\end{document}
